\documentclass[11pt]{article}
\usepackage[letterpaper,margin=1in]{geometry}
\usepackage[T1]{fontenc}
\usepackage{lmodern}
\usepackage{microtype}
\usepackage{setspace}
\usepackage{parskip}
\usepackage{amsmath,amssymb}
\usepackage{fancyhdr}
\usepackage{titlesec}
\usepackage{enumitem}
\usepackage{hyperref}

\hypersetup{hidelinks}
\setstretch{1.08}
\setlength{\parindent}{0pt}
\setlength{\parskip}{0.75em}
\pagestyle{fancy}
\fancyhf{}
\lhead{\small The Gedanbedai of Wake Island}
\rhead{\small Brian Tenneson}
\cfoot{\thepage}
\renewcommand{\headrulewidth}{0.3pt}
\titleformat{\section}{\large\bfseries}{}{0pt}{}

\begin{document}

\begin{titlepage}
\centering
\vspace*{1.4in}
{\Huge\bfseries The Gedanbedai of Wake Island\par}
\vspace{0.35in}
{\Large A Logic Parable\par}
\vspace{0.75in}
{\large Brian Tenneson\par}
\vspace{0.15in}
{\normalsize Story development and editorial assistance: ChatGPT (OpenAI)\par}
\vfill
{\itshape A story about language, self-description, authority, and what it means to say:\par}
\vspace{0.15in}
{\large ``I know that I know that I am a gedanbedai.''\par}
\vfill
September 2026
\end{titlepage}

\section*{The Gedanbedai of Wake Island}

Wake Island was not an island in the ordinary sense. It was a country so broad that a man could walk west for months and still be in Wake, and so thinly inhabited that two travelers meeting on the open road would sometimes stare at one another before remembering the proper greeting.

Millions lived there, but mostly in settlements separated by immense stretches of grass, stone, forest, and silence.

Leaving a settlement was regarded as foolish.

Usually it was.

But merchants left, because settlements needed one another. So did messengers, fugitives, surveyors, and---according to certain stories---the gedanbedai.

Every settlement had heard of the gedanbedai. Some had one. A few claimed several.

A gedanbedai was said to have studied storytelling for years. He could make an angry crowd laugh, a frightened crowd stand firm, or a victorious crowd become cautious. If the settlement required a historian, he remembered. If it required a prophet, he warned. If it required a fool, he made himself ridiculous.

No one knew how a gedanbedai was trained.

No one had witnessed an examination.

No one had attended a promotion.

Yet the gedanbedai spoke frequently about their training.

They spoke of an hour when words ceased to feel chosen and simply flowed through them, unimpeded.

One evening a trader named Terun arrived at a northern settlement after nineteen days alone on the road.

At supper he was seated beside an old woman named Mara.

``Are you the gedanbedai?'' Terun asked.

Mara smiled.

``I know that I am.''

Terun had heard that answer before.

So he tried the question that travelers were warned never to ask.

``How were you promoted?''

``I wasn't.''

Terun stared at her.

``Then you aren't a gedanbedai.''

``Why?''

``Because only someone promoted to gedanbedai is a gedanbedai.''

``Who told you that?''

``A gedanbedai.''

Mara laughed.

The next morning Terun found her speaking in the square. Two farmers had nearly started a riot over water rights. Within fifteen minutes Mara had persuaded each faction that allowing the other to speak first was a sign of superior courage.

The crowd dispersed peacefully.

Terun followed her.

``You certainly behave like a gedanbedai.''

``That is less interesting,'' said Mara, ``than whether I am one.''

``But surely somebody must decide.''

``Who?''

``The other gedanbedai.''

``Who promoted them?''

Terun was silent.

Mara continued walking.

After a while he said, ``Then perhaps there are no gedanbedai.''

``Perhaps.''

``But you just told me you knew that you were one.''

``I did.''

Terun stopped.

Mara walked three more steps before noticing.

``That is impossible,'' he said. ``If nobody has ever been promoted, and promotion is necessary, then you cannot know you are one.''

``Excellent,'' said Mara.

``Excellent?''

``You have discovered that one of our premises must be wrong.''

``Which one?''

``That,'' said Mara, ``is why the training takes years.''

Terun hurried after her.

``But wait. Do you know that you know you are a gedanbedai?''

Mara turned.

``Yes.''

Terun threw up his hands. ``That makes it worse!''

``Only if saying it makes it true.''

And she resumed walking.

Terun caught up again.

``So what checks it?''

``Nothing.''

``Nothing?''

``No examiner. No promotion ledger. No seal hidden in a temple. No official who can certify the sentence.''

``Then what does your second `I know' add?''

Mara looked pleased.

``At last,'' she said, ``you have asked a gedanbedai question.''

They walked together toward the edge of the settlement.

Beyond the final house stretched a road so empty that it seemed to divide the world in two.

Terun said, ``Then perhaps \emph{I know that I know that I am a gedanbedai} is not a certificate at all.''

``No.''

``It might merely be a story someone tells about himself.''

``Yes.''

``It might even be a very convincing story.''

``Yes.''

``And an entire settlement might believe it.''

``Certainly.''

``But none of that makes it knowledge.''

Mara stopped at the first milestone.

``You are learning quickly.''

Terun frowned. ``Does that mean I am becoming a gedanbedai?''

``No one becomes a gedanbedai.''

``You said that already.''

``Yes.''

``Then what happens?''

Mara looked back at the settlement.

``Eventually,'' she said, ``people begin asking you questions.''

And before Terun could ask another, she walked alone onto the road.

\newpage
\section*{A Short Logical Note}

The parable deliberately separates four things that ordinary language often merges:

\begin{enumerate}[leftmargin=1.5em]
\item behaving as though one has a status;
\item asserting that one has that status;
\item asserting that one knows one has that status; and
\item possessing a warrant that makes the knowledge claim factive under a declared epistemic rule.
\end{enumerate}

In the story, the sentence
\[
\text{``I know that I know that I am a gedanbedai''}
\]
may exist as language, conviction, social authority, or self-description even though the society supplies no external promotion record or verifier that by itself turns the assertion into certified knowledge.

The intended AMLD analogy is therefore not that self-description is impossible without a verifier. It is that the meaning of the word \emph{know} must be stated. A verifier-grounded AMLD treatment can require an accepted warrant for the inner claim and a further receipt or reflection warrant for the outer claim. Wake Island removes that machinery on purpose so that the reader is forced to ask the same question Terun finally asks:

\begin{center}
\emph{What does the second ``I know'' add?}
\end{center}

That is the riddle.

\end{document}
